% =============================================================================
% Appendix (simplified): length-filter ablations.
%
% Condensed version of appendix_length_filters.tex; labels are distinct, so both
% files can be included in the same document.
%
% Usage: \input this file after \appendix in the paper.
% Required packages: graphicx, booktabs, amsmath.
% Assumes UTF-8 input and T1 font encoding (both provided by the ACL template).
%
% Figures are generated by ablations/filtering/make_figures.py. By default they
% are read from ./figures; to change this, define the macro before the \input:
%   \newcommand{\LengthAblationFigDir}{ablations/filtering/figures}
% =============================================================================
\providecommand{\LengthAblationFigDir}{ablations/filtering/figures}

\section{Length Filtering Ablations}
\label{app:lensimple}

After structural validation, which also discards questions whose text or
choices contain an image, the pipeline applies three word-count filters,
in this order: the question must have between 4 and 1,000 words; the choices,
taken together, at most 300 words; and the question plus its choices at most
1,000 words (all limits inclusive). The lower limit removes fragments left by
extraction failures. The upper limits keep examples at a manageable size,
which matters when several are concatenated in few-shot prompts, without
discarding the long, passage-based questions that are typical of several
Brazilian exams. The thresholds were selected through three ablation studies
on the 41,987 structurally valid records (17 exams), summarized below. They
are coverage studies: they show what each threshold removes and from which
exams, not whether removal improves the benchmark.

\paragraph{How the thresholds were chosen.} Words were counted exactly as in
the pipeline, with NeMo Curator's \texttt{WordCountFilter}, which splits text
on whitespace. Each threshold was varied over a grid while we recorded the
share of records kept overall, by academic level, and for the most affected
exam, the last being the most informative number, since aggregate retention
stays close to 100\% near every candidate
(Figure~\ref{fig:lensimple-sweeps}, Table~\ref{tab:lensimple-candidates}).
For each length band around a candidate threshold, we also read eight
records, sampled with a fixed seed, to judge whether the items that would be
removed were defective or legitimate.

\begin{figure*}[t]
\centering
\includegraphics[width=\textwidth]{\LengthAblationFigDir/threshold_sweeps.pdf}
\caption{Share of records kept overall and for the most affected exam as each
threshold varies: (a)~minimum and (b)~maximum question length, (c)~maximum
choice length, and (d)~maximum combined length. $w_q$, $w_c$ and $w_t$ are
the word counts of the question, the choices and both; $D_0$ is the full set
of structurally valid records and $D_2$ the records left by the first two
filters. Dashed lines mark the selected thresholds.}
\label{fig:lensimple-sweeps}
\end{figure*}

\begin{table}[t]
\centering
\footnotesize
\setlength{\tabcolsep}{4pt}
\begin{tabular}{@{}rrrl@{}}
\toprule
Threshold & Removed & Kept (\%) & Most affected exam \\
\midrule
\multicolumn{4}{@{}l}{\emph{Question minimum}} \\
3 & 20 & 99.952 & COMVEST 99.349 \\
\textbf{4} & \textbf{105} & \textbf{99.750} & \textbf{COMVEST 98.698} \\
5 & 958 & 97.718 & OAB 91.859 \\
\midrule
\multicolumn{4}{@{}l}{\emph{Question maximum}} \\
750 & 409 & 99.026 & IME 82.305 \\
\textbf{1,000} & \textbf{152} & \textbf{99.638} & \textbf{IME 88.022} \\
1,500 & 49 & 99.883 & IME 95.554 \\
\midrule
\multicolumn{4}{@{}l}{\emph{Choices maximum}} \\
250 & 115 & 99.726 & CFCES 96.814 \\
\textbf{300} & \textbf{42} & \textbf{99.900} & \textbf{CFCES 97.922} \\
350 & 17 & 99.960 & CFCES 98.892 \\
\midrule
\multicolumn{4}{@{}l}{\emph{Question plus choices maximum}} \\
950 & 77 & 99.815 & AFA 96.992 \\
\textbf{1,000} & \textbf{36} & \textbf{99.914} & \textbf{IME 98.763} \\
1,050 & 9 & 99.978 & IME 99.794 \\
\bottomrule
\end{tabular}
\caption{Candidate thresholds around each selected value (bold). Each filter
is evaluated on its own over the 41,987 structurally valid records, except
the combined filter, which is evaluated on the 41,689 records left by the
first two. The last column gives the exam with the lowest share of records
kept (\%).}
\label{tab:lensimple-candidates}
\end{table}

\paragraph{Question length.} Questions with one to three words are
truncated stems, such as \emph{Assinale} or \emph{Imunidade:}, whereas
four-word questions are already complete instructions whose content lies in
the choices (\emph{Assinale a alternativa CORRETA:} alone occurs 141 times).
We therefore set the minimum to four words, which removes 105 records, raising
it to five would remove 958, mostly from OAB, whose retention would fall to
91.9\%. At the other end, the longest questions are legitimate reading
passages, mostly from the Portuguese section of IME, which repeats the same
text across items: the 152 questions over 1,000 words share only 29 passages.
The maximum of 1,000 words lies between the 99.5th and 99.9th percentiles
and is a size decision, not a quality judgment. Together, both limits remove
257 records, and IME is the most affected exam (88.0\% kept).

\paragraph{Choice length.} Very long sets of choices come from legitimately
verbose options (legal propositions, accounting rules, clinical plans), from
Markdown tables repeated in every option, where each cell separator counts as
a word, and occasionally from text of the next question leaking into an
option. The limit of 300 words sits at the 99.9th percentile. On its own it
removes 42 records, 41 of them beyond those already removed by the question
filter, and CFCES is the most affected exam (97.9\% kept). Because the tail mixes
legitimate options with errors, this limit is a conservative size cap rather
than an error detector.

\paragraph{Question plus choices.} The first two filters still admit
examples of up to 1,300 words. The last filter aligns the size of the whole
example with the 1,000-word ceiling used for the question. It removes 36
records, 12 each from AFA, BNDES and IME, all of them long reading passages.
As a consequence, in the length-filtered corpus this filter already implies the question
maximum, which is kept as a first-stage guard.

\begin{table}[t]
\centering
\footnotesize
\setlength{\tabcolsep}{4pt}
\begin{tabular}{@{}lrrr@{}}
\toprule
Stage & Removed & Retained & Kept (\%) \\
\midrule
Input & -- & 41,987 & 100.000 \\
Question (4--1,000) & 257 & 41,730 & 99.388 \\
Choices (${\le}\,300$) & 41 & 41,689 & 99.290 \\
Combined (${\le}\,1{,}000$) & 36 & 41,653 & 99.205 \\
\bottomrule
\end{tabular}
\caption{Records removed and kept after each length filter, applied in
pipeline order. \emph{Kept} is relative to the 41,987 input records.}
\label{tab:lensimple-stages}
\end{table}

\begin{figure}[t]
\centering
\includegraphics[width=\columnwidth]{\LengthAblationFigDir/removals_by_exam.pdf}
\caption{Records removed per exam, as a percentage of the exam, stacked by
the filter that removed them ($w_q$: question words; $w_c$: choice words;
$w_t$: both). Labels give removed and total records.}
\label{fig:lensimple-removals}
\end{figure}

\paragraph{Overall effect.} These counts describe the length-filtering
stage, before deduplication (Appendix~\ref{app:dedup-short}), item revision
(Appendix~\ref{app:rev-short}) and the construction of the published
splits. The three filters remove 334 records and keep
41,653 (99.205\%, see Table~\ref{tab:lensimple-stages}). They retain 98.1\% of the
high-school questions and 99.5\% of the undergraduate ones. The main side
effect falls on IME, which keeps 86.9\% of its questions because of its long
reading passages (Figure~\ref{fig:lensimple-removals}). Every other exam keeps
at least 97.4\%. Three exams (CNU, OBI and REVALIDA) lose no records. No
exam's share of the dataset changes by more than 0.33 percentage points.

\paragraph{Limitations.} Apart from the manual inspection, no quality labels
or downstream evaluations were used, so the studies cannot show that removal
improves the benchmark. Whitespace-separated words are only a proxy for model
tokens, and they overcount tables and formulas. The thresholds were chosen
after inspecting the results, favouring round values in the far tail. Finally,
because the filters act on individual questions, items that share a reading
passage can be split, with some kept and others removed.
